% !TeX root = ../main.tex
% Add the above to each chapter to make compiling the PDF easier in some editors.

\chapter{Introduction}\label{chapter:introduction}

Morphological systems encode how words are formed and how their forms vary with grammatical function. For languages with rich morphology, a single lemma may correspond to a large paradigm of inflected forms. Robust handling of morphology is therefore central to a wide range of downstream \ac{NLP} applications, including syntactic parsing, information extraction, and machine translation.

This thesis focuses on two core tasks that make morphological competence measurable: \emph{morphological inflection} and \emph{morphological analysis}. In the inflection task, the system receives a lemma together with a bundle of morphological features and must generate the corresponding inflected surface form (e.g., verb conjugation). In the analysis task, the system must recover the lemma and the \ac{MSD} from an observed inflected form, which thus can be seen as the inverse task to morphological inflection. In addition to the form alone, many modern settings provide sentence context, which can disambiguate morphological analyses that are underspecified at the token level.

The main goal of this work is to analyze how different modeling approaches perform on these two tasks across multiple typologically diverse languages, including both high- and low-resource settings, and to understand how data sources and annotation conventions affect the interpretation of results. Concretely, we evaluate non-neural and neural baselines alongside modern sequence-to-sequence models and prompting-based \ac{LLM} baselines.

\paragraph{Models.}
For morphological inflection, we compare (i) a non-neural SIGMORPHON baseline, (ii) a neural transducer baseline, and (iii) fine-tuned ByT5 models~\parencite{xue-etal-2022-byt5}. For morphological analysis, we evaluate ByT5 models trained in the inverse direction on UniMorph. Additionally, we perform context-based analysis experiments on \ac{UD}-derived data, including a context-aware ByT5 variant and inference-only LLM baselines. Finally, we include a bidirectional ByT5 model that is trained jointly on forward (inflection) and inverse (analysis) examples, enabling a unified investigation of multi-task learning for both mappings.

\paragraph{Datasets and evaluation.}
The experiments draw on two widely used morphological resources: UniMorph \parencite{batsuren-etal-2022-unimorph} and Universal Dependencies \parencite{nivre-etal-2020-universal}. Because these resources differ in annotation schemes and file formats, we construct unified train--development--test splits and a common instance representation for all model families. For the \ac{UD}-based setting, we map \ac{UD} feature bundles into a UniMorph-style tag inventory using an official conversion tool \parencite{mccarthy-etal-2018-marrying} when translators are available for the language. When translators are missing, \ac{UD} features cannot be converted into a comparable UniMorph-style \ac{MSD} inventory, which has direct implications for which metrics can be reported consistently across languages.

\paragraph{Contributions.}
This thesis makes the following contributions:
\begin{itemize}
	\item It provides a unified and reproducible experimental pipeline for training and evaluating a diverse set of models on morphological inflection and analysis across multiple languages.
	\item It offers a systematic comparison between established shared-task baselines, fine-tuned ByT5 models, and prompting-based LLM baselines under consistent preprocessing and evaluation.
	\item It highlights how dataset properties and annotation conventions, including differences between UniMorph and \ac{UD} and language-specific lemma conventions, can create apparent outliers and limit cross-resource comparability.
\end{itemize}

\paragraph{Thesis outline.}
Chapter~\ref{chapter:background} introduces the tasks and the baseline methods. Chapter~\ref{chapter:datasets} describes the datasets, languages, and preprocessing steps used to construct unified splits. Chapter~\ref{chapter:methodology} details the experimental setup, model configurations, and evaluation procedure. Chapter~\ref{chapter:results} presents and discusses the results for inflection and analysis. Finally, Chapter~\ref{chapter:conclusion} summarizes the main findings, discusses limitations, and outlines directions for future work.

The code used for the experiments is available at \url{https://github.com/muriloescher/bachelor-thesis}.